\chapter{Fundamentação Teórica e Trabalhos Relacionados}
\label{cap:conceitos}

Para compreender o impacto da Inteligência Artificial Generativa (IAG) no setor do entretenimento e justificar a necessidade da criação de uma taxonomia, é indispensável estabelecer uma base conceitual profunda. Este capítulo tem como objetivo nivelar o conhecimento do leitor sobre os dois universos que se cruzam nesta pesquisa: a arte cinematográfica e a ciência da computação. O texto apresenta as definições e as fases do ecossistema audiovisual, os fundamentos técnicos das arquiteturas generativas, as aplicações diretas dessas tecnologias no fluxo de trabalho e, por fim, explora o conceito de taxonomia e os trabalhos prévios da literatura.

\section{O Ecossistema Audiovisual}

O termo audiovisual designa qualquer meio de comunicação que combine simultaneamente elementos visuais e sonoros para transmitir uma mensagem, narrativa ou experiência estética \citep{chion2019audiovision}. Muito além de uma manifestação puramente artística, o audiovisual consolidou-se ao longo do último século como uma indústria complexa, que engloba o cinema, a televisão, a publicidade e, mais recentemente, as mídias digitais \citep{voinea2025}.

A criação de um produto audiovisual exige a orquestração de múltiplas disciplinas, incluindo a escrita, a fotografia, a atuação, o design de som e a engenharia de software. Para que esta pesquisa seja precisa em sua análise sobre a inserção de novas tecnologias, é necessário delimitar exatamente qual vertente do audiovisual está sendo estudada.

\subsection{Audiovisual Interativo versus Não Interativo}

Os produtos audiovisuais contemporâneos dividem-se em duas grandes categorias baseadas no nível de agência concedido ao espectador: os interativos e os não interativos \citep{lazarte2025}.

O audiovisual interativo é caracterizado pela capacidade do usuário de alterar ativamente o fluxo, a narrativa ou a perspectiva da obra em tempo real. O exemplo mais notório são os jogos digitais e as experiências de realidade virtual. Nesses formatos, o motor de jogo renderiza as imagens dinamicamente com base nas decisões do jogador (como mover a câmera ou escolher um caminho). A narrativa não possui uma duração temporal fixa e o desfecho pode variar.

Em contrapartida, o audiovisual não interativo, que é o objeto de estudo exclusivo desta pesquisa, refere-se a conteúdos de progressão linear \citep{gavran2025}. Filmes, séries de televisão, curtas-metragens, videoclipes e peças publicitárias enquadram-se nesta definição. Neste formato, o espectador assume uma postura passiva (receptiva). A agência criativa pertence inteiramente ao diretor e à equipe de produção, que definem previamente cada enquadramento, a duração exata de cada cena e a ordem dos eventos.

Esta restrição de linearidade impõe um rigor extremo à produção: como o espectador não pode controlar a câmera, qualquer erro de continuidade visual, falha de iluminação ou quebra na coerência temporal entre um quadro e o seguinte destrói a imersão. É justamente essa exigência por estabilidade temporal perfeita que torna a inserção da IA neste setor um desafio computacional tão complexo \citep{yin2025spatiotemporal}.

\subsection{O Pipeline de Produção Audiovisual Clássico}
\label{subsec:pipeline}

Para garantir que a obra final mantenha a coerência narrativa e estética do primeiro ao último minuto, a criação de conteúdos lineares não interativos segue um fluxo de trabalho (pipeline) fragmentado e altamente estruturado, tradicionalmente dividido em três fases principais \citep{orak2024, gavran2025}:

\textbf{Pré-produção}: É a fase fundamental de planejamento e ideação criativa. Tudo o que acontecerá na tela é decidido aqui. Esta etapa inicia-se com a escrita e o refinamento do roteiro. A partir do roteiro aprovado, artistas conceituais criam ilustrações (\textit{concept art}) para definir a identidade visual de cenários e personagens. Em seguida, os diretores criam o \textit{storyboard}, que é uma representação desenhada quadro a quadro de todas as cenas do filme, semelhante a uma história em quadrinhos. Muitas vezes, esses desenhos são animados com sons provisórios, criando um \textit{animatic}, que dita o ritmo exato que o filme terá antes mesmo de qualquer câmera ser ligada. Além disso, a pré-produção envolve o planejamento logístico, como a escolha de locações e testes de elenco.

\textbf{Produção}: É a fase de execução material e captação primária. Em filmes \textit{live-action} (com atores reais), esta é a etapa das filmagens no \textit{set}, envolvendo a direção de atores, o posicionamento físico de câmeras, captação de áudio direto e o design de iluminação. Em animações puras, é o momento em que os animadores constroem o movimento dos personagens em ambientes virtuais (2D ou 3D). É a fase mais cara e que envolve a maior quantidade de profissionais simultaneamente no mesmo espaço.

\textbf{Pós-produção}: É a fase de lapidação, montagem e finalização da obra, onde os fragmentos captados são costurados para formar a narrativa linear. Inicia-se com a edição de vídeo (montagem), onde o diretor escolhe as melhores tomadas e dita o ritmo final. Em paralelo, entram os Efeitos Visuais (VFX), que envolvem a criação de elementos por computador, a remoção de cabos e objetos indesejados, e a técnica de rotoscopia (o recorte minucioso de personagens do cenário original para inseri-los em novos fundos). A fase é concluída com a correção de cor (\textit{color grading}), que unifica a aparência de todas as câmeras, e o design de som, que inclui a adição de trilha sonora, efeitos sonoros (\textit{foley}) e eventuais dublagens de diálogos perdidos durante a filmagem.

\section{Fundamentos da Inteligência Artificial Generativa (IAG)}

Se o audiovisual tradicional baseia-se na captação do mundo físico ou na modelagem manual por artistas digitais, a IAG propõe um paradigma inteiramente novo de criação. A IAG é um subcampo do aprendizado de máquina (\textit{machine learning}) \citep{russell2016artificial}. Enquanto a IA clássica é analítica (treinada para detectar padrões, como reconhecer um rosto em uma foto), a IA Generativa é sintética: ela mapeia a estrutura estatística de milhões de dados pré-existentes e utiliza esse conhecimento para criar conteúdos inéditos (imagens, textos, vídeos) a partir do zero \citep{franganillo2022}.

A evolução dessas ferramentas no mercado audiovisual foi sustentada por três arquiteturas computacionais fundamentais:

\textbf{Redes Adversariais Generativas (GANs)}: Propostas em 2014, revolucionaram a síntese visual através de uma competição matemática. O modelo consiste em duas redes: um gerador, que tenta fabricar uma imagem falsa o mais realista possível, e um discriminador, que tenta estimar se a imagem é proveniente dos dados de treino ou se foi criada sinteticamente \citep{goodfellow2014generative}. Esse treinamento competitivo gera resultados impressionantes, mas as GANs são difíceis de controlar e frequentemente sofrem com instabilidade ou geram artefatos característicos durante o processo de aprendizagem, o que limita sua adoção para sequências complexas \citep{karras2020analyzing}.

\textbf{Transformers}: Ganhando evidência em 2017, os transformers destacam-se pela capacidade de compreender o contexto amplo de uma sequência de dados, baseando-se inteiramente em mecanismos de atenção para entender como as partes se relacionam entre si \citep{vaswani2017attention}. Inicialmente desenvolvidos para textos, eles permitiram que as máquinas passassem a entender instruções complexas em linguagem natural e realizassem tarefas com poucos exemplos (\textit{few-shot learning}) \citep{brown2020language}. São a base que permite aos usuários interagirem com as ferramentas atuais apenas descrevendo o que desejam em formato de texto (\textit{prompts}) \citep{floridi2020gpt3}.

\textbf{Modelos de Difusão}: Consolidados na síntese de alta resolução, tornaram-se o estado da arte absoluto para geração visual \citep{dhariwal2021diffusion, rombach2022high}. Eles operam por um princípio de destruição e reconstrução: o processo de formação da imagem é decomposto na aplicação sequencial de redes que removem o ruído de uma estática abstrata \citep{rombach2022high}. A rede neural aprende então a reverter esse processo passo a passo. Esta arquitetura garante uma qualidade visual superior às GANs e permite um controle refinado através de mecanismos de orientação (\textit{guidance}), permitindo direcionar estilos artísticos e composições sem a necessidade de retreinar o modelo \citep{dhariwal2021diffusion}.

Foi a popularização destas arquiteturas, somada à facilidade das interfaces textuais introduzidas no final de 2022, que tirou a IA Generativa dos laboratórios acadêmicos e a colocou diretamente nos estúdios de produção e nas ilhas de edição \citep{dwivedi2023chatgpt}.

Foi nesse contexto de amadurecimento técnico que ocorreu um ponto de inflexão decisivo para a popularização da IAG. O lançamento do \textit{ChatGPT}, desenvolvido pela OpenAI e disponibilizado publicamente no final de 2022, tornou acessível ao grande público a interação direta com modelos de linguagem de larga escala baseados em arquiteturas Transformer \citep{brown2020language}. Esse marco não apenas ampliou a adoção das tecnologias generativas, mas também consolidou um novo paradigma de interação humano-máquina centrado em \textit{prompts}, influenciando diretamente o desenvolvimento de ferramentas audiovisuais baseadas em texto para imagem e texto para vídeo \citep{dwivedi2023chatgpt}.

\section{Aplicações da IAG no Pipeline Audiovisual}

Com o amadurecimento dos Modelos de Difusão e dos Transformers, a adoção de ferramentas generativas deixou de ser um experimento e passou a infiltrar-se ativamente no fluxo de trabalho tradicional do audiovisual não interativo. O impacto da IAG, no entanto, não ocorre de maneira uniforme; ele varia drasticamente dependendo da fase da produção.

A fase de planejamento é, atualmente, a que mais absorve as tecnologias generativas com sucesso \citep{voinea2025}. Como a pré-produção foca na experimentação, pequenas inconsistências visuais da IA são não apenas toleradas, como muitas vezes bem-vindas para inspirar novas ideias.
Modelos de linguagem são amplamente utilizados para auxiliar roteiristas na superação de bloqueios criativos, estruturação de arcos narrativos e formatação de textos \citep{dwivedi2023chatgpt}. Simultaneamente, ferramentas de geração de imagens estão a ser integradas nos departamentos de arte. Diretores utilizam \textit{prompts} textuais para gerar \textit{concept arts} de cenários, figurinos e paletas de cores em questão de segundos. Além disso, plataformas de geração visual estão automatizando a criação de \textit{storyboards}, substituindo horas de ilustração manual por rascunhos sintéticos que ditam a composição de câmeras que será filmada posteriormente.

A fase de filmagem é a que sofre menos interferência direta das tecnologias puramente generativas, uma vez que depende da presença física de luz, atores e cenários. Sua principal interseção com a IAG está na renderização neural, que permite reiluminar cenas computacionalmente após a captação, sem a necessidade de remontar fisicamente o esquema de iluminação no \textit{set} \citep{einabadi2021relighting}.

A pós-produção concentra a maior promessa de automação e, paradoxalmente, os maiores desafios técnicos \citep{gavran2025}. A IAG tem sido empregada com êxito como ferramenta de assistência pontual: algoritmos aumentam a resolução de imagens antigas, removem automaticamente objetos indesejados do cenário e limpam ruídos de áudio captados nas filmagens de forma muito superior às ferramentas digitais clássicas.
A rotoscopia, uma das tarefas mais manuais e exaustivas do cinema, começa a ser acelerada por modelos de aprendizagem profunda capazes de identificar e rastrear contornos de sujeitos complexos automaticamente. Na área do som, modelos generativos de áudio conseguem clonar a voz dos atores, permitindo corrigir palavras mal pronunciadas no \textit{set} sem a necessidade de reconvocar o elenco para regravações em estúdio de dublagem. Contudo, a automação total, como a geração autônoma de longas cenas de vídeo apenas com comandos de texto, ainda encontra limitações severas nesta fase \citep{yin2025spatiotemporal}.

\section{Desafios Técnicos, Operacionais e Éticos}

Apesar das evidências de adoção crescente e das projeções otimistas sobre os benefícios do uso por indivíduos e pelo mercado \citep{grandview2025generative, gambacorta2025exploring}, a inserção da IAG no audiovisual linear está longe de ser um processo sem atritos. A academia e o mercado enfrentam desafios que ultrapassam a esfera da engenharia de software, exigindo análises multidimensionais \citep{dwivedi2023chatgpt}.

O principal obstáculo técnico é a falta de coerência temporal \citep{gavran2025, orak2024}. Ferramentas baseadas em difusão lutam para manter a consistência da identidade de um objeto ao longo do tempo em um vídeo. Uma camisa listrada gerada por IA pode mudar o número de listras de um \textit{frame} para o outro, um fenômeno conhecido na indústria como "cintilação" ou \textit{flickering}. Para o audiovisual não interativo, onde a suspensão de descrença do espectador é frágil, essa instabilidade restringe o uso de vídeos 100% gerados por IA a peças muito curtas, experimentais ou publicitárias.

Do ponto de vista operacional, as interfaces atuais dependem de \textit{prompts} textuais baseados em modelos de linguagem \citep{brown2020language}. A linguagem humana é ambígua, o que gera imprevisibilidade nos resultados da máquina. Profissionais de pós-produção necessitam de ferramentas precisas e determinísticas; a natureza probabilística da IAG dificulta o seu encaixe em cronogramas de estúdio com prazos rígidos \citep{franganillo2022}.

Por fim, os desafios éticos representam o maior gargalo para a consolidação institucional destas tecnologias. A esmagadora maioria dos modelos é treinada absorvendo o estilo fotográfico, de animação e de direção de artistas humanos reais sem que estes tenham consentido ou sido remunerados. O uso comercial dessas ferramentas em estúdios expõe as produções a riscos de plágio, indefinições sobre direitos autorais e falta de transparência sobre a origem sintética do material, desafios evidenciados pelas pesquisas na área de identificação de \textit{deepfakes} \citep{floridi2020gpt3, franganillo2022, tiwari2023}. Paralelamente, a facilitação extrema de processos laboriosos coloca em risco a sustentabilidade de departamentos criativos, gerando preocupações severas sobre o aumento de desigualdades e a substituição de mão de obra especializada em domínios de trabalho intensivo em informação \citep{capraro2024impact, dwivedi2023chatgpt}.

\section{O Papel da Taxonomia e a Lacuna na Literatura}

Diante de um ecossistema com elevada complexidade técnica, diversidade de aplicações ao longo do pipeline audiovisual e implicações éticas significativas, torna-se necessária a adoção de mecanismos estruturados de organização do conhecimento. Nesse contexto, a taxonomia configura-se como uma abordagem fundamental para sistematizar domínios emergentes.

Uma taxonomia pode ser definida como um sistema formal de classificação que organiza entidades em dimensões e categorias estruturadas, permitindo a análise comparativa e a identificação de relações entre os elementos \citep{Wieringa2006}. Diferentemente de abordagens puramente descritivas, como revisões sistemáticas da literatura, a taxonomia estabelece critérios explícitos de categorização e busca estruturar o domínio de forma consistente.

Nos últimos anos, diferentes taxonomias têm sido propostas para organizar aspectos específicos da Inteligência Artificial e, mais recentemente, da Inteligência Artificial Genertiva. Por exemplo, \citet{strobel2024} propõem uma taxonomia de aplicações de IAG baseada em seus papéis funcionais, categorizando sistemas como geradores, assistentes e sintetizadores. Essa abordagem contribui para a compreensão do papel da tecnologia no ecossistema digital, porém não considera o contexto operacional específico de domínios como o audiovisual.

Paralelamente, estudos voltados à governança e segurança da IA têm desenvolvido taxonomias focadas em riscos e impactos. Trabalhos como os de \citet{slattery2024} organizam riscos de sistemas de IA segundo dimensões como causa, domínio e intencionalidade, enquanto \citet{marchal2024} propõem classificações para usos indevidos de modelos generativos, estruturando padrões de abuso e exploração da tecnologia. Essas abordagens são relevantes para a compreensão dos impactos sociotécnicos, mas não estabelecem conexão com os aspectos técnicos ou operacionais da produção audiovisual.

No domínio específico da produção audiovisual, a literatura apresenta esforços de organização do conhecimento principalmente por meio de revisões sistemáticas, como os trabalhos de \citet{lazarte2025} e \citet{orak2024}, que categorizam aplicações da IAG em diferentes etapas do pipeline ou tipos de uso. Contudo, essas classificações emergem de forma descritiva e não configuram taxonomias formais, uma vez que não definem dimensões estruturadas nem critérios explícitos de organização.

Dessa forma, observa-se que as abordagens existentes são predominantemente fragmentadas, focando isoladamente em dimensões funcionais, técnicas ou éticas. Não foi identificada, até o momento, uma taxonomia que integre de maneira sistemática os aspectos técnicos (arquiteturas e capacidades), operacionais (etapas do pipeline audiovisual) e sociotécnicos (riscos, impactos e papel humano) da adoção da IAG no audiovisual não interativo.

Essa lacuna estrutural limita a capacidade de análise comparativa entre soluções, dificulta a compreensão das trocas envolvidas na adoção da tecnologia e compromete a construção de diretrizes mais robustas para seu uso. Nesse contexto, justifica-se a proposição de uma taxonomia multidimensional que articule essas diferentes dimensões de forma integrada, constituindo a principal contribuição deste trabalho: a TIGA, apresentada no Capítulo~\ref{cap:taxonomia}.